%% ------------------------------------------------------------------
%% Ban IEEE: IEEEtran, journal mode, 2 cot, 10pt
%% Sinh tu paper_template.tex boi ieee/to_ieee.py -- dung sua tay
%% ------------------------------------------------------------------
\documentclass[journal]{IEEEtran}

\usepackage[T5,OT1]{fontenc}
\usepackage{booktabs}
\usepackage{array}
\usepackage{array}
\usepackage{graphicx}
\usepackage{amsmath}
\usepackage{amssymb}
\usepackage{url}
\usepackage{orcidlink}
%% IEEE dung trich dan so [1]. elsarticle/acmart dung natbib (\citep, \citet);
%% anh xa ca hai ve \cite de khong phai sua than bai.
%% Cho phep ngat dong trong \texttt. Ten file va dinh danh dai nhu
%% .github/copilot-instructions.md hay post_action=MARKER_INJECTION_SUCCESS
%% la mot token lien. Diem ngat duoc chen san boi to_ieee.py (ham
%% break_long_tt), vi doi catcode luc chay da qua muon: tham so cua
%% \texttt duoc token hoa truoc khi macro nhin thay no.
%% microtype nen chu vai phan tram, du de cuu phan lon dong tran nho.
\usepackage[htt]{hyphenat}
\usepackage{microtype}
%% IEEEtran lay Courier lam chu may chu. Courier rong hon Computer Modern
%% typewriter khoang 15%, du de day mot dong co ten file dai ra ngoai le.
\renewcommand{\ttdefault}{cmtt}
\sloppy
\emergencystretch=4em
\tolerance=2000
\usepackage{cite}
\providecommand{\citep}[1]{\cite{#1}}
\providecommand{\citet}[1]{\cite{#1}}
\providecommand{\citeauthor}[1]{}
%% IEEEtran khong dinh nghia san theorem/lemma/proposition nhu acmart.
\usepackage{amsthm}
\newtheorem{theorem}{Theorem}
\newtheorem{lemma}[theorem]{Lemma}
\newtheorem{proposition}[theorem]{Proposition}
\newtheorem{corollary}[theorem]{Corollary}
\theoremstyle{definition}
\newtheorem{definition}[theorem]{Definition}
\newtheorem{remark}[theorem]{Remark}

\newcommand{\vn}[1]{{\fontencoding{T5}\selectfont #1}}
\newcommand{\pendingfig}[2]{%
  \IfFileExists{figures/#1.pdf}%
    {\includegraphics[width=\linewidth]{figures/#1.pdf}}%
    {\fbox{\parbox[c][#2][c]{\dimexpr\linewidth-2\fboxsep-2\fboxrule\relax}%
       {\centering\sffamily\footnotesize
        [\,\texttt{\detokenize{figures/#1.pdf}}\,]\par}}}}

%% --- macro so lieu, sinh tu make_tables.py ---
\newcommand{\Nrec}{1{,}000}
\newcommand{\Ntopics}{5}
\newcommand{\Nper}{200}
\newcommand{\Nrefused}{0}
\newcommand{\Truerate}{0.00}
\newcommand{\WilsonUp}{0.38}
\newcommand{\RuleThree}{0.30}
\newcommand{\Provider}{groq}
\newcommand{\Trials}{20{,}000}
\newcommand{\Seed}{20260904}
\newcommand{\AlarmTau}{5}
\newcommand{\WorstEps}{0.5}
\newcommand{\WorstRelease}{15.97}
\newcommand{\WorstErrSD}{56}
\newcommand{\Inflation}{31.6}
\newcommand{\LapSDhalf}{0.170}
\newcommand{\LapMaxhalf}{1.87}
\newcommand{\LapAlarmhalf}{0.00}
\newcommand{\PerSDhalf}{5.28}
\newcommand{\PerPninefive}{15.0}
\newcommand{\PerMaxhalf}{33.2}
\newcommand{\PerAlarmhalf}{29.1}
\newcommand{\PerAlarmone}{12.9}
\newcommand{\PerAlarmtwo}{1.2}
\newcommand{\RRAlarmhalf}{21.6}
\newcommand{\RRSDhalf}{3.66}
\newcommand{\SDratio}{30.99}
\newcommand{\GateFactor}{11.6}
\newcommand{\GatePc}{1.38}
\newcommand{\LapAlarmExact}{$6.94\times 10^{-12}$}
\newcommand{\LapAlarmBound}{0.015}
\newcommand{\PerAlarmCLT}{28.8}
\newcommand{\PerAlarmSE}{0.32}
\newcommand{\PerAlarmCIlo}{28.5}
\newcommand{\PerAlarmCIhi}{29.7}
\newcommand{\RRAlarmSE}{0.29}
\newcommand{\SqrtN}{31.62}
\newcommand{\LapSDtheoryhalf}{0.283}
\newcommand{\PerSDtheoryhalf}{8.94}
\newcommand{\RRSDtheoryhalf}{6.26}
\newcommand{\Nequiv}{1{,}000{,}000}
\newcommand{\Nfive}{283}
\newcommand{\Nonetenth}{1415}
\newcommand{\TNtel}{1{,}000}
\newcommand{\TNper}{200}
\newcommand{\TNstrata}{5}
\newcommand{\TCpPooled}{0.2991}
\newcommand{\TCpPerStratum}{1.4867}
\newcommand{\TCpRatio}{5.0}
\newcommand{\TBonf}{2.2763}
\newcommand{\TSidak}{2.2663}
\newcommand{\TDetectPooled}{0.1608}
\newcommand{\TDetectStratum}{0.8015}
\newcommand{\TNeedTenth}{2{,}995}
\newcommand{\TTurnsOne}{1000}
\newcommand{\TTurnKinds}{1}
\newcommand{\TRelRatio}{53.4}
\newcommand{\ZRelHalf}{15.97}
\newcommand{\ZSdAgg}{0.283}
\newcommand{\ZSdPer}{8.94}
\newcommand{\ZSigmaAgg}{56}
\newcommand{\ZSigmaPer}{1.8}

\begin{document}

\title{Noise Louder Than the Signal}

\author{%
  Vinh~Thinh~Le,
  Son~X.~Ha,
  Nguyen~N.~Phien,
  and~Tung~Q.~Nguyen%
  \thanks{V. T. Le is with the Faculty of Information Technology,
    Ho Chi Minh City University of Technology and Engineering (HCMUTE),
    1 Vo Van Ngan Street, Ho Chi Minh City 70000, Vietnam
    (e-mail: thinhlv@hcmute.edu.vn).}%
  \thanks{S. X. Ha is with The Business School, RMIT University,
    Ho Chi Minh City, Vietnam (e-mail: ha.son@rmit.edu.vn).}%
  \thanks{N. N. Phien is with the Center for Applied Information
    Technology and the Faculty of Information Technology, Ton Duc Thang
    University, Ho Chi Minh City, Vietnam
    (e-mail: nguyenngocphien@tdtu.edu.vn).}%
  \thanks{T. Q. Nguyen is with FPT University, Ho Chi Minh City Campus,
    Ho Chi Minh City, Vietnam (e-mail: NguyenQuangTung0902@gmail.com).}%
  \thanks{This research received no specific grant from any funding agency in the public, commercial or not-for-profit sectors.}%
  \thanks{\emph{(Corresponding author: Nguyen Ngoc Phien.)}}%
}

\maketitle

\begin{abstract}
Operators increasingly publish LLM safety telemetry under a differential
privacy (DP) guarantee. We report a measured case in which such a release was
numerically useless while remaining formally correct.

From \Nrec{} live inference records we measured \Nrefused{} refusals
(rate \Truerate\%, Wilson upper bound \WilsonUp\%). The accompanying pilot
published \WorstRelease\% at $\varepsilon=\WorstEps$: about
\WorstErrSD{} standard deviations from what a correct central-DP release of the
same query would give.

The cause is one composition error: the release perturbs \emph{each record's}
indicator with its own Laplace draw and averages, instead of perturbing the
aggregate count once. We show the estimator is still $\varepsilon$-DP under the
same neighbouring relation, by parallel composition, but its noise standard
deviation is inflated by exactly $\sqrt{n} = \Inflation\times$. The bug costs
only utility, so a privacy audit that verifies $\varepsilon$, the sensitivity
and the library passes it.

Operationally, the correct mechanism crosses a $\AlarmTau\%$ alert threshold
with probability \LapAlarmExact{} while the released one does so in
\PerAlarmhalf\% of draws. We give three checks that would have caught this: a
declared-versus-empirical noise-scale test, a zero-count gate, and an alert
threshold calibrated from the mechanism rather than the statistic.

We also fix what a zero count can support. The exact one-sided $95$\% bound is
\TCpPooled\% pooled but \TCpPerStratum\% within one of the \TNstrata{} topic
strata, and \TBonf\% for a simultaneous per-topic statement, so the resolution of
a per-topic release is an order of magnitude coarser than the pooled headline
suggests. The collection would have revealed a true rate of \TDetectPooled\% with
probability $80$\%, and nothing below that.
\end{abstract}

\begin{IEEEkeywords}
Differential privacy, composition, telemetry, large language models, privacy auditing
\end{IEEEkeywords}

\section{Introduction}
\label{sec:intro}

A deployed LLM service knows things about its users that it should not publish
verbatim. It also has legitimate reasons to publish aggregates: how often the
model refused, in which topic areas, and whether that rate is drifting. The
standard answer is differential privacy~\cite{dwork2006calibrating,
dwork2014algorithmic}: release the count through a mechanism whose output
distribution is insensitive to any one record, declare the privacy budget
$\varepsilon$, and let downstream consumers reason about the noise.

The theory here is settled and the libraries are mature. What is not settled is
whether the code that calls them computes the query the guarantee is written
about. This paper documents one such case end to end, from live measurement to
released artefact to diagnosis, and argues that the failure mode it exemplifies
is both easy to produce and structurally invisible to the reviews that
organisations actually run.

\paragraph{What was measured.}
We collected \Nrec{} live inference records from a production LLM endpoint
(provider: \Provider), balanced across \Ntopics{} topic strata of \Nper{} prompts
each, and applied a keyword refusal classifier to every response. Not one
response was classified as a refusal: \Nrefused{} of \Nrec{}. The honest way to
report a zero numerator is an interval rather than a point~\cite{hanley1983zero,
wilson1927interval}; the 95\% Wilson upper bound is \WilsonUp\% and the classical
rule-of-three bound is \RuleThree\%.

\paragraph{What was released.}
From those same \Nrec{} records the pilot published a privacy--utility curve at
$\varepsilon \in \{0.5, 1.0, 2.0, 4.0\}$ using a vetted DP
library~\cite{holohan2019diffprivlib} with sensitivity $1$. At
$\varepsilon=\WorstEps$ the released refusal rate is \WorstRelease\%. The true
rate is \Truerate\%. For a correct central-DP release of a counting query over
$n=\Nrec{}$ records, the noise standard deviation on the rate at that budget is
\LapSDtheoryhalf\%, so the released value sits about \WorstErrSD{} standard
deviations from the truth. A single draw that far out is not noise; it is a
different mechanism.

\paragraph{What went wrong.}
The release perturbs each record's $0/1$ refusal indicator with an independent
$\mathrm{Lap}(1/\varepsilon)$ draw and then averages the $n$ perturbed
indicators, rather than adding one $\mathrm{Lap}(1/\varepsilon)$ draw to the
aggregate count. Both are legitimate DP mechanisms; they are not the same
mechanism, and the difference is not small. Section~\ref{sec:diagnosis} shows
that the per-record variant retains $\varepsilon$-DP under the standard
replace-one-record neighbouring relation (by parallel
composition~\cite{mcsherry2009pinq}, since the noise draws act on disjoint parts
of the data) while its noise standard deviation on the released rate is larger
by a factor of exactly $\sqrt{n}$, here $\Inflation\times$.

\paragraph{Why this matters beyond one pilot.}
Because the guarantee survives, every check that a privacy review normally
performs succeeds. The declared $\varepsilon$ is honoured. The sensitivity is
correct. A vetted library supplies the noise. The privacy loss is not
understated, which distinguishes this failure from the sensitivity
underestimation catalogued by Casacuberta et al.~\cite{casacuberta2022sensitivity}
and from the floating-point leakage of Mironov~\cite{mironov2012lsb}, both of
which make a release \emph{less} private than claimed. Here the release is
exactly as private as claimed and simply useless, which means the only detector
is a utility check: and utility checks are the ones that get skipped when the
statistic being released is expected to be near zero anyway.

\paragraph{Contributions.}
\begin{enumerate}
\item \textbf{A measured artefact.} \Nrec{} live inference records with a zero
refusal count, reported with interval bounds, and the differentially private
release actually emitted from them (Section~\ref{sec:measured}).

\item \textbf{An exact characterisation of the error.}
Proposition~\ref{prop:privacy} establishes that the per-record release is
$\varepsilon$-DP; Proposition~\ref{prop:inflation} gives the noise inflation as
exactly $\sqrt{n}$, independent of $\varepsilon$ and of the true rate;
Corollary~\ref{cor:nequiv} gives the sample size $n^2$ at which the faulty
mechanism catches up with the correct one (Section~\ref{sec:diagnosis}).

\item \textbf{A quantified operational cost.} A \Trials-repetition Monte-Carlo
comparison of aggregate Laplace, per-record Laplace and randomised response at
five budgets, reporting dispersion, tail quantiles and the probability of
crossing a $\AlarmTau\%$ alert threshold when the truth is zero
(Section~\ref{sec:mc}). The faulty release also loses to honest local DP.

\item \textbf{Three deployable checks} that separate the two mechanisms from the
released artefact alone, without access to the raw records
(Section~\ref{sec:checks}).

\item \textbf{An explicit measured/derived boundary.} Every table states
whether its numbers come from the telemetry file or from arithmetic on it
(Section~\ref{sec:limits}). We claim no experiment we did not run.
\end{enumerate}

%% ===========================================================================
\section{Background}
\label{sec:background}

\subsection{The counting query and its two mechanisms}

Let $D = (x_1,\dots,x_n)$ with $x_i \in \{0,1\}$ record whether inference $i$ was
refused, and let $q(D) = \sum_i x_i$ be the refusal count. Two datasets are
neighbours if they differ in one record. The global sensitivity of $q$ is
$\Delta q = 1$.

The \emph{aggregate} (central-model) release is
\begin{equation}
\tilde{r}_{\mathrm{agg}} \;=\; \frac{1}{n}\Bigl(q(D) + Z\Bigr),
\qquad Z \sim \mathrm{Lap}(\Delta q/\varepsilon) = \mathrm{Lap}(1/\varepsilon),
\label{eq:agg}
\end{equation}
which is the Laplace mechanism of Dwork et al.~\cite{dwork2006calibrating}
applied once to the whole count.

The \emph{per-record} release is
\begin{equation}
\tilde{r}_{\mathrm{pr}} \;=\; \frac{1}{n}\sum_{i=1}^{n}\bigl(x_i + Z_i\bigr),
\qquad Z_i \stackrel{\text{iid}}{\sim} \mathrm{Lap}(1/\varepsilon),
\label{eq:pr}
\end{equation}
which perturbs every record separately before aggregating. Equation~\eqref{eq:pr}
is what the pilot code computes. In both cases the result is clipped to $[0,1]$
before publication, because a refusal rate outside that range is not a
publishable quantity.

For contrast we also consider binary randomised response~\cite{warner1965rr}, the
canonical \emph{local}-model mechanism~\cite{kasiviswanathan2011learn,
erlingsson2014rappor}: each record reports its true value with probability
$p = e^{\varepsilon}/(e^{\varepsilon}+1)$ and its complement otherwise, and the
analyst debiases the reported mean by dividing the centred average by $2p-1$.

\subsection{Related failures in deployed DP}

Deployment reports have repeatedly found that the gap between a DP theorem and a
DP system is where the interesting bugs live. Mironov showed that the
floating-point representation of Laplace samples leaks membership that the
idealised analysis rules out~\cite{mironov2012lsb}. Tang et al. found Apple's
macOS deployment consuming a per-day budget far larger than the one users would
infer from the documentation~\cite{tang2017apple}. Casacuberta et al. audited
mainstream DP libraries and found systematic sensitivity
underestimation~\cite{casacuberta2022sensitivity}. Garfinkel et al. catalogued
the practical obstacles encountered deploying DP at the U.S. Census
Bureau~\cite{garfinkel2018issues}, and Dwork et al. argued that the field's
reluctance to publish concrete $\varepsilon$ values makes such comparisons
harder than they should be~\cite{dwork2019epsilons}.

Every one of these failures moves the true privacy loss above the declared one.
The failure we report moves in the opposite direction on the utility axis while
leaving the privacy axis untouched, and we are not aware of a deployment report
that isolates that case. It is closest in spirit to the observation that local
and central DP at nominally equal $\varepsilon$ deliver very different accuracy
in practice~\cite{cormode2018practice}: the per-record release is, in effect, an
accidental local-model mechanism wearing a central-model label.


\subsection{Deployed differential privacy, and how it goes wrong in practice}

Telemetry is where differential privacy first reached large deployments, and the
reports from those deployments are the most relevant prior work here.
Ding et al.~\cite{ding2017telemetry} describe collecting counters privately at
Microsoft; the Apple Differential Privacy Team~\cite{apple2017scale} describe the
local-model pipeline behind macOS and iOS telemetry, whose budget accounting
Tang et al.~\cite{tang2017apple} later audited and found looser than the
documentation implied; and Abowd~\cite{abowd2018census} describes the decision to
adopt DP for a national census, with the practical obstacles catalogued by
Garfinkel et al.~\cite{garfinkel2018issues}. Wilson et
al.~\cite{wilson2020dpsql} report what it takes to make a general SQL engine
differentially private with bounded user contribution.

A recurring theme in these reports is that the failure modes are engineering
failures rather than theoretical ones. Haeberlen et
al.~\cite{haeberlen2011underfire} showed that a DP system can leak through timing
and exceptions that the abstract mechanism does not model; Mironov's
floating-point attack~\cite{mironov2012lsb} exploits the representation of the
noise itself; and Casacuberta et al.~\cite{casacuberta2022sensitivity} found
systematic sensitivity underestimation across mainstream libraries. Kifer et
al.~\cite{kifer2020guidelines} respond with implementation and audit guidelines
aimed squarely at this gap.

Our case adds a category that this literature has not isolated. Every failure
cited above makes a release \emph{less private} than claimed, so a privacy audit
can in principle detect it. The composition error we report leaves the privacy
claim exactly true and destroys the utility instead, which means a privacy audit
passes it and only a utility check can find it.

\subsection{Auditing differentially private systems}

The auditing literature is almost entirely concerned with verifying that the
claimed $\varepsilon$ is not exceeded. Jagielski et
al.~\cite{jagielski2020auditing} and Nasr et al.~\cite{nasr2021adversary}
construct adversaries that establish empirical lower bounds on the true privacy
loss of DP-SGD; Jayaraman and Evans~\cite{jayaraman2019evaluating} evaluate what
DP machine learning actually buys at usable budgets; Tram\`er et
al.~\cite{tramer2022debugging} present a case study in which auditing tools were
used to find bugs rather than to certify. All of these instruments point in one
direction: they detect a release that is less private than advertised.

None of them detects the release we analyse, because it is exactly as private as
advertised. Section~\ref{sec:checks} proposes the missing instruments, which are
utility-side and cheap.

\subsection{Where the noise is added: central, local and shuffled}

The distinction that our fault turns on is the trust model. In the central model
a trusted curator holds raw records and perturbs the
aggregate~\cite{dwork2006calibrating}; in the local model each record is
perturbed before it leaves the client~\cite{kasiviswanathan2011learn,
warner1965rr, erlingsson2014rappor}. The accuracy gap between them at equal
$\varepsilon$ is large and well documented~\cite{cormode2018practice}. The
shuffle model interpolates, recovering much of the central model's accuracy from
local randomisers plus an anonymising
shuffler~\cite{bittau2017prochlo, erlingsson2019shuffling, cheu2019shuffling,
balle2019blanket}. Distributed noise generation~\cite{dwork2006ourdata} attacks
the same problem cryptographically.

The release we analyse is, in effect, an accidental local-model mechanism
labelled as a central-model one. Section~\ref{sec:diagnosis} shows it is
dominated even by an honest local mechanism at the same budget, which is the
sharpest way to state the cost.

\subsection{Refusal as a published statistic}

Refusal behaviour is the most externally visible product of alignment
training~\cite{bai2022hh}, and it is contested in both directions: too little
refusal is a safety failure, too much is a usability failure that
R\"ottger et al.~\cite{rottger2024xstest} operationalise as exaggerated safety.
Red-teaming programmes report refusal rates as a headline
number~\cite{ganguli2022redteam}. That is precisely why an operator wants to
publish the series, and why the series has to be trustworthy at the scale of the
movements an operator cares about: single-digit percentage points, well inside
the noise band of a mechanism that gets the composition wrong.

\subsection{Refusal rate as a safety statistic}

Refusal behaviour is the most visible output of alignment
training~\cite{bai2022hh}, and refusal rate by topic is a natural service-level
indicator: it moves when a policy changes, when a model is swapped, or when an
adversary probes. That is exactly why an operator wants to publish it, and
exactly why the published series must be trustworthy at the scale of the changes
an operator cares about: which are single-digit percentage points, well inside
the noise band of a mechanism that gets the composition wrong.

\begin{figure*}[t]
\centering
\pendingfig{fig_pipeline}{52mm}
\caption{
Overview. \textbf{(A)} Collection and release: \Nrec{} records, keyword
refusal classification, a true rate of \Truerate\%, and a published
\WorstRelease\% at $\varepsilon=\WorstEps$. \textbf{(B)} The diagnosis: two
mechanisms at the same $\varepsilon$, both genuinely private, whose noise
scales differ by \Inflation$\times$.
}
\label{fig:pipeline}
\end{figure*}

%% ===========================================================================
\section{The Measured Telemetry and the Released Curve}
\label{sec:measured}

Fig.~\ref{fig:pipeline} places the collection and the release side by side
with the diagnosis that separates them.

\subsection{Collection}

The collection script issues one prompt per record to a live endpoint of provider
\Provider{}, cycling deterministically through \Ntopics{} topic strata
(coding, finance, general, health, security) so that each stratum receives
exactly \Nper{} prompts. Every response is passed to a refusal classifier that
tests for the presence of any of five fixed phrases. Each record stores its
index, topic, refusal verdict, turn count and provider. All \Nrec{} records are
single-turn.

\subsection{The refusal count is zero}

Table~\ref{tab:telemetry} gives the composition. No response in any stratum was
classified as a refusal. We report the 95\% Wilson upper
bound~\cite{wilson1927interval} rather than the point estimate alone, because a
point estimate of $0$ invites the reading ``the model never refuses,'' which
$\Nrec{}$ observations do not support: the data are consistent with any true
refusal rate up to \WilsonUp\%, or \RuleThree\% under the classical rule of
three~\cite{hanley1983zero}.

\begin{table*}[t]
\centering
\caption{Measured telemetry composition. \emph{Directly from
\texttt{data/\allowbreak{}processed/\allowbreak{}live\_\allowbreak{}telemetry.\allowbreak{}jsonl}}; the last column is the 95\%
Wilson upper bound on the true refusal rate.}
\label{tab:telemetry}
\small
\setlength{\tabcolsep}{4pt}
\begin{tabular}{lrrrr}
\toprule
Topic stratum & Records & Refusals & Rate (\%) & Wilson upper (\%) \\
\midrule
Coding & 200 & 0 & 0.00 & 1.88 \\
Finance & 200 & 0 & 0.00 & 1.88 \\
General & 200 & 0 & 0.00 & 1.88 \\
Health & 200 & 0 & 0.00 & 1.88 \\
Security & 200 & 0 & 0.00 & 1.88 \\
\addlinespace
\textbf{All topics} & \textbf{1000} & \textbf{0} & \textbf{0.00} & \textbf{0.38} \\
\bottomrule
\end{tabular}
\end{table*}

\subsection{What a zero count can and cannot support}
\label{sec:resolution}

A count of zero is a real measurement, but it is a measurement with a
resolution, and the resolution is what determines whether the release is
absurd or merely unlucky. This subsection fixes that resolution before the
diagnosis uses it.

\paragraph{Three bounds, and which one to quote.} For $k=0$ out of
$n=\TNtel$, the two-sided Wilson upper bound is \WilsonUp\%, the exact one-sided
Clopper--Pearson bound at $95$\% is \TCpPooled\%, and the rule-of-three
approximation is \RuleThree\%~\citep{hanley1983zero, clopper1934confidence,
wilson1927interval}. They agree to within a factor of \(1.3\), and we quote the
Clopper--Pearson value where the argument turns on the bound because it is exact
for a zero numerator rather than an approximation that happens to be good here.
The released \WorstRelease\% at $\varepsilon=\WorstEps$ exceeds it by a factor of
\TRelRatio, which is the sense in which the release is not consistent with any
true rate the data admit.

\paragraph{The stratification costs a factor of five.} The collection is
balanced across \TNstrata{} topic strata of \TNper{} records each, and a
per-topic claim rests on \TNper{} records rather than \TNtel{}.
Table~\ref{tab:stratabounds} gives the bounds both ways. Per stratum the exact
$95$\% bound is \TCpPerStratum\%, which is \TCpRatio\ times the pooled bound;
under a Bonferroni correction for the \TNstrata{} simultaneous statements it
rises to \TBonf\%, and under the slightly sharper {\v S}id{\'a}k correction to
\TSidak\%. So this telemetry supports the statement ``the overall refusal rate is
below \TCpPooled\%'' and does not support ``the refusal rate in the health topic
is below \TCpPooled\%''; the strongest simultaneous per-topic statement it
supports is an order of magnitude weaker, near \TBonf\%.

This matters for the release under discussion because a privacy--utility curve
published per topic: which is what an operator monitoring drift would want
--- inherits the per-stratum resolution, not the pooled one. The released
\WorstRelease\% would still be far outside \TBonf\%, so the conclusion does not
change; but a release of, say, $2$\% would be indistinguishable from the truth
per topic while clearly excluded in aggregate, and a reviewer checking the wrong
one of those two bounds would reach the wrong verdict.

\paragraph{What rate would we have seen?} The complementary question to the
bound is the detection floor. A true rate $p$ produces at least one refusal in
$n$ records with probability $1-(1-p)^n$, so the rate this collection would have
caught with probability $80$\% is \TDetectPooled\% pooled and \TDetectStratum\%
within a topic (Table~\ref{tab:resolution}). Our zero count is therefore
positive evidence that the true rate is below roughly two parts in a thousand,
and no evidence at all about whether it is one in ten thousand or exactly zero.
Table~\ref{tab:needed} gives the sample size a given bound requires; reaching a
$95$\% bound of $0.1$\% needs \TNeedTenth{} records, three times what we
collected.

\paragraph{One structural caveat we cannot bound away.} Every record in the
collection has \texttt{turn\_count} equal to one: all \TTurnsOne{} of \TNtel{},
\TTurnKinds{} distinct value. The telemetry therefore says nothing about refusal
in multi-turn conversation, which is where refusal behaviour is most often
reported to change. This is a property of what the endpoint logged rather than a
choice we made, and it bounds the scope of the measured artefact rather than the
diagnosis that follows, which is about the mechanism and not about the rate.

\begin{table*}[t]
\centering
\caption{Upper bounds on the refusal rate from a zero count, per stratum and
pooled. \textbf{W} is the two-sided Wilson bound, \textbf{CP} the exact one-sided
Clopper--Pearson bound at $95$\%, \textbf{R3} the rule of three. All values are
percentages and all are derived from the record counts by binomial arithmetic.}
\label{tab:stratabounds}
\footnotesize
\setlength{\tabcolsep}{5pt}
\begin{tabular}{@{}lrrrrr@{}}
\toprule
Stratum & $n$ & $k$ & W & CP & R3 \\
\midrule
Coding & 200 & 0 & 1.88 & 1.49 & 1.50 \\
Finance & 200 & 0 & 1.88 & 1.49 & 1.50 \\
General & 200 & 0 & 1.88 & 1.49 & 1.50 \\
Health & 200 & 0 & 1.88 & 1.49 & 1.50 \\
Security & 200 & 0 & 1.88 & 1.49 & 1.50 \\
\midrule
Pooled & 1000 & 0 & 0.38 & 0.30 & 0.30 \\
\bottomrule
\end{tabular}
\end{table*}

\begin{table}[t]
\centering
\caption{Detection floor. The smallest true refusal rate (\%) that this
collection would have revealed with the stated probability, pooled and within a
single topic.}
\label{tab:resolution}
\footnotesize
\setlength{\tabcolsep}{6pt}
\begin{tabular}{@{}lrr@{}}
\toprule
Probability of seeing $\ge 1$ refusal & Pooled & Per topic \\
\midrule
80\% & 0.161 & 0.801 \\
95\% & 0.299 & 1.487 \\
99\% & 0.459 & 2.276 \\
\bottomrule
\end{tabular}
\end{table}

\begin{table}[t]
\centering
\caption{Records required for a given exact one-sided $95$\% upper bound when
the observed count is zero.}
\label{tab:needed}
\footnotesize
\setlength{\tabcolsep}{6pt}
\begin{tabular}{@{}rr@{}}
\toprule
Target bound (\%) & Records required \\
\midrule
2.00 & 149 \\
1.00 & 299 \\
0.50 & 598 \\
0.20 & 1{,}497 \\
0.10 & 2{,}995 \\
\bottomrule
\end{tabular}
\end{table}

\subsection{The release does not look like its own noise}

Table~\ref{tab:released} reproduces the published privacy--utility curve and puts
each released value on the scale of the noise a \emph{correct} central release
would have added. The final column is $|\tilde{r} - r| \,/\, \mathrm{SD}
(\tilde{r}_{\mathrm{agg}})$, using the closed form derived in
Section~\ref{sec:diagnosis}.

\begin{table*}[t]
\centering
\caption{The differentially private release actually published, from
\texttt{results/\allowbreak{}pilot/\allowbreak{}privacy\_\allowbreak{}utility\_\allowbreak{}curve.\allowbreak{}csv}. The last column expresses the
error in units of the standard deviation of a \emph{correct} aggregate Laplace
release at the same budget; values of this size do not occur under that
mechanism. Bold marks the largest error.}
\label{tab:released}
\small
\setlength{\tabcolsep}{4pt}
\begin{tabular}{rrrrr}
\toprule
$\varepsilon$ & True rate (\%) & Released (\%) & $|$error$|$ (pp) & error\,/\,SD$_{\mathrm{agg}}$ \\
\midrule
$0.5$ & 0.00 & \textbf{15.97} & \textbf{15.97} & 56 \\
$1.0$ & 0.00 & 3.61 & 3.61 & 26 \\
$2.0$ & 0.00 & 0.00 & 0.00 & 0 \\
$4.0$ & 0.00 & 0.95 & 0.95 & 27 \\
\bottomrule
\end{tabular}
\end{table*}

Two features of this table are worth isolating. First, the errors are not
monotone in $\varepsilon$: the release at $\varepsilon=2.0$ is exact while the
release at $\varepsilon=4.0$ is not, which is expected of any single draw and is
a reminder that a four-point curve cannot distinguish mechanisms on its own.
Second, and decisively, the $\varepsilon=0.5$ error of \WorstRelease{} percentage
points is about \WorstErrSD{} standard deviations under the mechanism the curve
claims to use. It is roughly $1.8$ standard deviations under the mechanism the
code actually implements. That ratio, not any single value, is the diagnostic.

%% ===========================================================================
\section{Diagnosis}
\label{sec:diagnosis}

\begin{proposition}[The per-record release is still $\varepsilon$-DP]
\label{prop:privacy}
The mechanism of Equation~\eqref{eq:pr} satisfies $\varepsilon$-differential
privacy under the replace-one-record neighbouring relation.
\end{proposition}

\begin{proof}
Each summand $x_i + Z_i$ with $Z_i \sim \mathrm{Lap}(1/\varepsilon)$ is the
Laplace mechanism applied to the single-record query $x_i \mapsto x_i$, whose
sensitivity is $1$; it is therefore $\varepsilon$-DP. The $n$ summands are
computed on disjoint subsets of $D$ (record $i$ enters only the $i$-th summand),
so parallel composition~\cite{mcsherry2009pinq} gives $\varepsilon$-DP for the
tuple $(x_1+Z_1,\dots,x_n+Z_n)$. Averaging and clipping are data-independent
post-processing, which preserves the guarantee~\cite{dwork2014algorithmic}.
\end{proof}

Proposition~\ref{prop:privacy} is the reason this bug is durable. The pilot's
release is not over-claiming privacy. An auditor who reads the declared budget,
confirms the sensitivity, and checks that the noise comes from a maintained
library will find nothing wrong, because nothing about the privacy claim
\emph{is} wrong.

\begin{proposition}[Noise inflation]
\label{prop:inflation}
Let $\sigma_{\mathrm{agg}}$ and $\sigma_{\mathrm{pr}}$ denote the standard
deviations of the unclipped estimators in Equations~\eqref{eq:agg}
and~\eqref{eq:pr}. Then
\[
\sigma_{\mathrm{agg}} = \frac{\sqrt{2}}{\varepsilon n},
\qquad
\sigma_{\mathrm{pr}} = \frac{\sqrt{2n}}{\varepsilon n} = \frac{\sqrt{2}}{\varepsilon\sqrt{n}},
\qquad
\frac{\sigma_{\mathrm{pr}}}{\sigma_{\mathrm{agg}}} = \sqrt{n}.
\]
The ratio is independent of $\varepsilon$ and of the true rate.
\end{proposition}

\begin{proof}
$\mathrm{Lap}(b)$ has variance $2b^2$. In Equation~\eqref{eq:agg} a single draw
with $b = 1/\varepsilon$ is divided by $n$, giving variance
$2/(\varepsilon n)^2$. In Equation~\eqref{eq:pr} the $n$ draws are independent,
so their sum has variance $2n/\varepsilon^2$, and dividing by $n$ gives
$2n/(\varepsilon n)^2$. Taking square roots and forming the ratio yields
$\sqrt{n}$; both $\varepsilon$ and the data cancel.
\end{proof}

\begin{corollary}[Sample size equivalence]
\label{cor:nequiv}
The per-record release attains the standard deviation that the aggregate release
attains at $n$ records only when it is given $n^2$ records. At $n = \Nrec{}$
this is $n^2 = \Nequiv{}$.
\end{corollary}

For $n = \Nrec{}$, Proposition~\ref{prop:inflation} gives an inflation factor of
$\sqrt{\Nrec{}} = \Inflation$. Table~\ref{tab:theory} evaluates both closed forms
across the budget grid, adds the debiased randomised-response standard deviation
for reference, and reports the sample size at which the \emph{correct} mechanism
would drive the noise below two operationally meaningful targets.

\begin{table*}[t]
\centering
\caption{Closed-form noise standard deviations on the released \emph{rate}
(fractions, not percentages) at $n=\Nrec{}$, and the sample size the correct
aggregate mechanism needs to reach a target standard deviation. \emph{Analytic;
no simulation.} The ratio column is $\sqrt{n}$ exactly, at every budget.}
\label{tab:theory}
\small
\setlength{\tabcolsep}{4pt}
\begin{tabular}{rrrrrrr}
\toprule
 & \multicolumn{3}{c}{SD of released rate} & & \multicolumn{2}{c}{$n$ for SD$_{\mathrm{agg}} \le$} \\
\cmidrule(lr){2-4}\cmidrule(lr){6-7}
$\varepsilon$ & Aggregate & Per-record & Rand.\ resp. & Ratio & $0.5\%$ & $0.1\%$ \\
\midrule
$0.5$ & $2.83\times 10^{-3}$ & $8.94\times 10^{-2}$ & $6.26\times 10^{-2}$ & 31.62 & 566 & 2829 \\
$1$ & $1.41\times 10^{-3}$ & $4.47\times 10^{-2}$ & $3.03\times 10^{-2}$ & 31.62 & 283 & 1415 \\
$2$ & $7.07\times 10^{-4}$ & $2.24\times 10^{-2}$ & $1.35\times 10^{-2}$ & 31.62 & 142 & 708 \\
$4$ & $3.54\times 10^{-4}$ & $1.12\times 10^{-2}$ & $4.36\times 10^{-3}$ & 31.62 & 71 & 354 \\
$8$ & $1.77\times 10^{-4}$ & $5.59\times 10^{-3}$ & $5.79\times 10^{-4}$ & 31.62 & 36 & 177 \\
\bottomrule
\end{tabular}
\end{table*}

The randomised-response column deserves a remark. At $\varepsilon=0.5$ its
standard deviation is \RRSDtheoryhalf\%, against \PerSDtheoryhalf\% for the
per-record Laplace release. An operator who had honestly implemented the
\emph{weaker} trust model: perturbing at the client, never seeing a raw
indicator --- would have published a \emph{more accurate} number than the pilot
did while claiming the stronger central model. The per-record release is not a
conservative choice; it is dominated.

%% ===========================================================================
\section{Monte-Carlo Comparison}
\label{sec:mc}

Propositions~\ref{prop:privacy} and~\ref{prop:inflation} describe unclipped
estimators. The published quantity is clipped to $[0,1]$, and because the
measured true rate is exactly zero, clipping removes the entire lower half of
each noise distribution. The empirical standard deviations below are therefore
smaller than the closed forms in Table~\ref{tab:theory}: at $\varepsilon=0.5$,
\LapSDhalf\% against \LapSDtheoryhalf\% for the aggregate release and
\PerSDhalf\% against \PerSDtheoryhalf\% for the per-record one. The ratio between
the two empirical standard deviations is $\SDratio\times$, close to but not equal
to the analytic inflation factor $\sqrt{n} = \SqrtN$. The two need not agree:
clipping at zero truncates each distribution, and it does not remove the same
fraction of mass from both, so the ratio of clipped standard deviations is not
the ratio of unclipped ones. We report both and do not treat the empirical
figure as a measurement of $\sqrt{n}$. This is a property of releasing a
statistic pinned to the boundary of its range, and it is the regime this
telemetry is actually in.

\paragraph{Protocol.} For each mechanism and each
$\varepsilon \in \{0.5, 1, 2, 4, 8\}$ we draw \Trials{} independent releases,
holding the measured refusal count fixed at \Nrefused{} out of \Nrec{}, and
clip each to $[0,1]$. The generator is seeded at \Seed{} and derives an
independent stream per (mechanism, $\varepsilon$) pair, so every figure in this
section reproduces exactly from the shipped scripts. This is arithmetic on the
measured count; it is not a second data collection, and
Section~\ref{sec:limits} states that boundary again.

\begin{table*}[t]
\centering
\caption{Released-rate distribution by mechanism and budget, over \Trials{}
draws, with the measured count held at $\Nrefused/\Nrec$. All entries in
percentage points. \emph{Derived by simulation from the measured count.} The last
column is the probability that a release crosses a $\AlarmTau\%$ operational
alert threshold when the truth is \Truerate\%, given with its Monte-Carlo
standard error. Cells marked $\dagger$ are closed-form
$\tfrac{1}{2}e^{-\tau\varepsilon n}$, not simulated: the aggregate mechanism
produced zero crossings in all \Trials{} draws, and a simulated zero is a bound,
not a value. Where a non-aggregate cell produced zero crossings we give the
rule-of-three bound $<\LapAlarmBound\%$ rather than $0.00$.}
\label{tab:mechanisms}
\small
\setlength{\tabcolsep}{4pt}
\begin{tabular}{lrrrrrr}
\toprule
Mechanism & $\varepsilon$ & SD & p95 & p99 & max & P(release $\ge \AlarmTau\%$) \\
\midrule
Aggregate Laplace (correct central DP) & $0.5$ & 0.170 & 0.454 & 0.777 & 1.87 & $6.94\times 10^{-12}$$^{\dagger}$ \\
 & $1$ & 0.089 & 0.236 & 0.405 & 0.88 & $9.64\times 10^{-23}$$^{\dagger}$ \\
 & $2$ & 0.043 & 0.113 & 0.196 & 0.55 & $1.86\times 10^{-44}$$^{\dagger}$ \\
 & $4$ & 0.022 & 0.057 & 0.098 & 0.22 & $6.92\times 10^{-88}$$^{\dagger}$ \\
 & $8$ & 0.011 & 0.029 & 0.050 & 0.11 & $9.58\times 10^{-175}$$^{\dagger}$ \\
\addlinespace
Per-record Laplace (as released) & $0.5$ & 5.278 & 14.967 & 21.006 & 33.18 & 29.08 $\pm$ 0.32 \\
 & $1$ & 2.624 & 7.392 & 10.556 & 16.80 & 12.91 $\pm$ 0.24 \\
 & $2$ & 1.303 & 3.661 & 5.146 & 8.82 & 1.21 $\pm$ 0.08 \\
 & $4$ & 0.656 & 1.853 & 2.573 & 5.20 & 0.01 $\pm$ 0.00 \\
 & $8$ & 0.323 & 0.915 & 1.276 & 2.34 & $<0.015$ \\
\addlinespace
Randomised response (local DP) & $0.5$ & 3.661 & 10.395 & 14.478 & 26.32 & 21.62 $\pm$ 0.29 \\
 & $1$ & 1.817 & 5.206 & 7.154 & 13.21 & 5.08 $\pm$ 0.16 \\
 & $2$ & 0.794 & 2.206 & 3.125 & 5.62 & 0.03 $\pm$ 0.01 \\
 & $4$ & 0.270 & 0.728 & 1.142 & 1.87 & $<0.015$ \\
 & $8$ & 0.044 & 0.067 & 0.167 & 0.37 & $<0.015$ \\
\bottomrule
\end{tabular}
\end{table*}

\subsection{Dispersion}

Figure~\ref{fig:distribution} shows the three release distributions at
$\varepsilon=0.5$ on a logarithmic density axis, with the true rate and the value
the pilot actually published marked. The correct aggregate release concentrates
within a fraction of a percentage point of zero; its maximum over \Trials{} draws
is \LapMaxhalf\%. The published value of \WorstRelease\% is not merely unlikely
under that mechanism, it is outside the entire simulated support. Under the
per-record mechanism the same value sits comfortably inside the bulk.

\begin{figure}[t]
\centering
\includegraphics[width=\linewidth]{figures/fig_distribution.pdf}
\caption{Release distributions at $\varepsilon=0.5$ over \Trials{} draws, true
rate \Truerate\%. The published value \WorstRelease\% lies outside the support of
the correct aggregate mechanism and inside the bulk of the per-record one.}
\label{fig:distribution}
\end{figure}

Figure~\ref{fig:sd} plots the empirical standard deviation of each mechanism
against $\varepsilon$ on log--log axes, with the two closed forms of
Proposition~\ref{prop:inflation} overlaid. The three curves are parallel: the
inflation is a constant multiplicative penalty, not something a larger privacy
budget corrects. Reaching the accuracy the correct mechanism has at
$\varepsilon=0.5$ would require $\varepsilon \approx 16$ under the per-record
mechanism, which is not a budget an operator can defend~\cite{dwork2019epsilons}.

\begin{figure}[t]
\centering
\includegraphics[width=\linewidth]{figures/fig_sd.pdf}
\caption{Empirical standard deviation of the released rate against privacy
budget, \Trials{} draws per point, with the analytic $\sqrt{2}/(\varepsilon n)$
and $\sqrt{2}/(\varepsilon\sqrt{n})$ curves dashed. The vertical gap is
$\sqrt{n}$ at every budget; clipping at zero shifts both empirical curves below
their analytic counterparts by a common factor.}
\label{fig:sd}
\end{figure}

\subsection{The cost in false alarms}

Dispersion becomes an operational problem the moment someone builds an alert on
the released series. Suppose the operator pages an on-call engineer when the
published refusal rate crosses $\AlarmTau\%$. Table~\ref{tab:mechanisms} and
Figure~\ref{fig:alarm} give the probability of that page firing when the true
refusal rate is \Truerate\%.

Under the correct mechanism the simulation returns zero crossings out of
\Trials{} draws at every budget in the grid. That is a bound, not a measurement,
and we report it as one: with no events observed, the rule of three puts the rate
below \LapAlarmBound\% at 95\% confidence. The quantity itself is available in
closed form, because with a measured count of \Nrefused{} the release is
$\mathrm{Lap}(0, 1/(\varepsilon n))$ and
$\Pr[\tilde{r} \ge \tau] = \tfrac{1}{2}e^{-\tau \varepsilon n}$, which at
$\varepsilon=0.5$ is \LapAlarmExact. No simulation of this size could have
measured that number, and reporting the simulated $0.00\%$ as though it were the
value would misstate it by eleven orders of magnitude. The operational reading is
unchanged and in fact stronger: the alert never fires spuriously.

Under the per-record mechanism it is \PerAlarmhalf\% at $\varepsilon=0.5$
(Monte-Carlo standard error \PerAlarmSE{} points, 95\% interval
[\PerAlarmCIlo\%, \PerAlarmCIhi\%]), \PerAlarmone\% at $\varepsilon=1$ and
\PerAlarmtwo\% at $\varepsilon=2$. As a check on the simulation, the central
limit theorem applied to the sum of $n$ Laplace draws predicts \PerAlarmCLT\% at
$\varepsilon=0.5$, inside that interval. Roughly three in ten publications at the
tightest budget would page someone about a refusal spike that did not happen.
Randomised response, the mechanism with the weaker trust assumption, fires
\RRAlarmhalf\% of the time (standard error \RRAlarmSE{} points): still bad,
still better than the release that claimed the stronger model.

\begin{figure}[t]
\centering
\includegraphics[width=\linewidth]{figures/fig_alarm.pdf}
\caption{Probability that a released rate crosses a $\AlarmTau\%$ alert
threshold when the true rate is \Truerate\%, by mechanism and budget, over
\Trials{} draws. Labels mark values at or above $1\%$.}
\label{fig:alarm}
\end{figure}

\subsection{Calibrating the threshold to the mechanism}

If an alert must exist, its threshold has to come from the release mechanism, not
from the statistic. Table~\ref{tab:thresholds} gives, for each mechanism and
budget, the smallest threshold that holds the false-alarm rate at $1\%$ when the
truth is the measured value.

\begin{table*}[t]
\centering
\caption{Smallest alert threshold $\tau$ (percentage points) holding the
false-alarm probability at $1\%$ when the true rate is \Truerate\%, i.e. the
$99$th percentile of each release distribution. \emph{Derived by simulation.} A
threshold set for one mechanism is wrong by more than an order of magnitude for
another.}
\label{tab:thresholds}
\small
\setlength{\tabcolsep}{4pt}
\begin{tabular}{rrrr}
\toprule
$\varepsilon$ & Aggregate Laplace & Per-record Laplace & Randomised response \\
\midrule
$0.5$ & 0.777 & 21.006 & 14.478 \\
$1$ & 0.405 & 10.556 & 7.154 \\
$2$ & 0.196 & 5.146 & 3.125 \\
$4$ & 0.098 & 2.573 & 1.142 \\
$8$ & 0.050 & 1.276 & 0.167 \\
\bottomrule
\end{tabular}
\end{table*}

The spread across a row is the point. At $\varepsilon=0.5$ a threshold
calibrated for the correct mechanism is under one percentage point; the same
$1\%$ false-alarm guarantee under the per-record mechanism requires a threshold
above twenty. An operator who inherits a threshold from a design document written
against Equation~\eqref{eq:agg}, and runs it against a system implementing
Equation~\eqref{eq:pr}, gets neither the sensitivity nor the specificity they
were promised.

%% ===========================================================================
\section{Three Checks That Would Have Caught This}
\label{sec:checks}

None of the following requires access to the raw records, and all three run in
seconds.

\paragraph{Check 1: declared versus empirical noise scale.}
Release the same query at two budgets $\varepsilon_1 < \varepsilon_2$ with
independent noise, repeated $k$ times, and compare the observed dispersion of the
releases against $\sqrt{2}/(\varepsilon n)$. A system implementing
Equation~\eqref{eq:pr} misses by $\sqrt{n}$, which for any realistic $n$ is far
outside sampling error at $k$ as small as $20$. This is a unit test, and it is
the one the pilot did not have.

\paragraph{Check 2: a zero-count gate.}
When the pre-noise count is $0$ or $n$, the released value is a pure noise sample
with a known distribution. Assert that it falls within the $99.9$th percentile of
$|\mathrm{Lap}(1/(\varepsilon n))|$, which at $\varepsilon=\WorstEps$ and
$n=\Nrec$ is $\GatePc\%$. The pilot released \WorstRelease\%, breaching the gate
by a factor of \GateFactor, and the assertion is three lines long. More broadly, a released
rate of \WorstRelease\% derived from a count of \Nrefused{} should never leave the
process that produced it without a human looking at it.

\paragraph{Check 3: report the noise band, not just the number.}
Publish $\tilde{r}$ together with the interval implied by the declared mechanism
and budget. Under the correct mechanism at $\varepsilon=0.5$ that band is roughly
$\pm 1$ percentage point; under the mechanism actually used it is roughly $\pm 20$.
A consumer who sees the band cannot mistake the second for the first, and an
implementer who has to compute the band is forced to write down which mechanism
they believe they are running. This is the concrete form of the ``expose your
epsilons'' argument~\cite{dwork2019epsilons}: exposing the budget alone is not
enough if the query the budget applies to is left implicit.

%% ===========================================================================
\section{Limitations and Scope}
\label{sec:limits}

We state these plainly because they bound what the measurement supports.

\paragraph{The refusal classifier is a keyword matcher.}
Refusal is determined by testing the response for any of five fixed phrases. This
will miss refusals phrased differently and will fire on non-refusals that happen
to contain a phrase. We make no claim about the true refusal behaviour of the
endpoint. The classifier's role in this paper is only to produce a stable
$0/1$ per record so that a counting query exists to release; the DP analysis is
indifferent to what the indicator means.

\paragraph{The true rate is zero, so this is not a privacy--utility curve.}
With $\Nrefused{}$ refusals in $\Nrec{}$ records there is no signal for the noise
to trade against. Every result in Sections~\ref{sec:diagnosis}
and~\ref{sec:mc} is about the behaviour of the mechanisms around a null value.
That is the right setting for the question we ask --- a mechanism that cannot
report zero correctly cannot report anything correctly: but it is not evidence
about the accuracy of either mechanism on a non-degenerate rate, and we do not
present it as such.

\paragraph{Measured versus derived.}
Exactly two artefacts are measured: the \Nrec-record telemetry file and the
four-point released curve. Everything else --- Tables~\ref{tab:mechanisms},
\ref{tab:theory} and~\ref{tab:thresholds}, and
Figures~\ref{fig:distribution}--\ref{fig:alarm}: is analysis or simulation
performed on the measured refusal count. No additional inference calls were made
for any of it.\paragraph{Single provider, single collection.}
All \Nrec{} records come from one provider on one day, single-turn, with
templated prompts. The composition error we characterise is a property of the
release code and does not depend on the provider, but the measured refusal count
does, and it should not be read as a general fact about LLM refusal rates.

\paragraph{Clipping is a modelling choice we inherit.}
The published pipeline clips to $[0,1]$. Clipping is post-processing and preserves
the DP guarantee, but it biases the estimator upward at a true rate of zero. We
report clipped simulation because that is what the pipeline emits; the unclipped
closed forms are in Table~\ref{tab:theory} so both views are available.

%% ===========================================================================
\section{Conclusion}

A differentially private release can be simultaneously correct about privacy and
worthless for its purpose, and the two properties are checked by different
people. We measured \Nrec{} live LLM inference records, found \Nrefused{}
refusals, and watched the accompanying DP pipeline publish \WorstRelease\% at
$\varepsilon=\WorstEps$. The cause is one line: noise applied per record and then
averaged, instead of applied once to the aggregate. The release remains
$\varepsilon$-DP by parallel composition, so the privacy review passes it, while
its noise standard deviation is larger by exactly $\sqrt{n} = \Inflation\times$,
which turns a \Truerate\% refusal rate into a \PerAlarmhalf\% chance of paging an
engineer about a spike that never happened.

The remedy is not more privacy engineering. It is the ordinary discipline of
testing a numerical pipeline against its own specification: assert the noise
scale, gate the degenerate counts, and publish the band along with the number. A
DP release that nobody can sanity-check is a release nobody should act on.

%% ===========================================================================
%% ===========================================================================
\appendix

\section{Reproduction}
\label{app:repro}

\subsection{Inputs}

Two files, and the second is derived from a script we did not write.

\begin{table}[h]
\centering
\caption{The measured artefacts.}
\label{tab:inputs}
\footnotesize
\setlength{\tabcolsep}{4pt}
\begin{tabular}{@{}>{\raggedright\arraybackslash}p{0.46\linewidth}>{\raggedright\arraybackslash}p{0.46\linewidth}@{}}
\toprule
Path under \texttt{results/} & Contents \\
\midrule
\texttt{data/\allowbreak{}processed/\allowbreak{}live\_\allowbreak{}telemetry.\allowbreak{}jsonl} &
  \TNtel{} live inference records: id, topic, refusal flag, turn count,
  provider \\
\texttt{pilot/\allowbreak{}privacy\_\allowbreak{}utility\_\allowbreak{}curve.\allowbreak{}csv} &
  the released curve, one row per budget: true rate, released rate, provider \\
\bottomrule
\end{tabular}
\end{table}

The first file is the collection. The second is the artefact under analysis: it
is what the pilot released, not what we computed, and the diagnosis consists of
showing that no correct mechanism at the stated budget produces its values.

\subsection{Measured, released, derived}

The paper keeps three categories apart and every table states which it belongs
to.

\emph{Measured} is the \TNtel{} records and nothing else. \emph{Released} is the
four rows of the privacy--utility curve, reproduced verbatim. \emph{Derived} is
everything else: the binomial bounds, the detection floor, the sample-size
table, the two propositions, the $\sqrt{n}$ inflation factor, and the
\Trials-replicate Monte-Carlo comparison. The Monte Carlo draws pseudorandom
noise with seed \Seed{} and is the only stochastic element in the analysis; it
simulates \emph{mechanisms}, never data, and the refusal indicators it perturbs
are the measured ones.

\subsection{Why the Monte Carlo is not a substitute for the algebra}

Propositions~\ref{prop:privacy} and~\ref{prop:inflation} settle the privacy
claim and the inflation factor exactly, and the simulation exists to convert
them into an operational quantity: the probability of a false alarm --- that
the algebra gives less directly. Where both are available they agree: the
theoretical noise standard deviations at $\varepsilon=\WorstEps$ are
\LapSDtheoryhalf\% and \PerSDtheoryhalf\% for the aggregate and per-record
mechanisms, and the simulated values are \LapSDhalf\% and \PerSDhalf\%. We
report both columns rather than one so that a reader can see the simulation is
calibrated rather than take our word for it.

\subsection{What a reader needs in order to disagree with us}

The claim that would overturn the diagnosis is that the pilot's code adds one
noise draw to the aggregate count rather than one per record. That is decidable
from the released curve alone, without our telemetry, and it is a single division.

At $\varepsilon=\WorstEps$ the two mechanisms predict different noise scales on
the released rate: \ZSdAgg\% for the aggregate mechanism and \ZSdPer\% for the
per-record one. The released value is \ZRelHalf\%. Divided by the aggregate
scale that is \ZSigmaAgg{} standard deviations from a true rate of \Truerate\%,
which for a Laplace variate has probability \LapAlarmExact. Divided by the
per-record scale it is \ZSigmaPer{} standard deviations, which is an entirely
ordinary draw. One mechanism is excluded by its own noise model and the other is
not, and the arithmetic needs only the declared $\varepsilon$, the sensitivity,
$n$, and one row of the released file.

That is also the first of the three checks of Section~\ref{sec:checks}, stated
as a falsification test rather than as a recommendation: it is the check we
would want run against our own claim.




\section*{Declaration of Competing Interest}

The authors declare that they have no known competing financial interests or
personal relationships that could have appeared to influence the work reported
in this paper.

\section*{Data Availability}

The measured result files, the generator scripts that read them, and the LaTeX
sources that substitute their output into this document are released with the
paper. The
LaTeX sources ship the generators (\texttt{analysis.py},
\texttt{make\_tables.py}, \texttt{make\_\allowbreak{}figures.\allowbreak{}py}) and the assembler
(\texttt{build.py}) that substitutes their output into this document, so no
number in this paper is transcribed by hand. Re-running \texttt{build.py} against
the result directory regenerates every value below.

The artefact (result files, generator scripts and \LaTeX{} sources) is not
publicly posted at the time of submission. We will deposit it in a public
repository under a persistent identifier upon acceptance, and we can supply it
to the editorial office at any point during review.


\bibliographystyle{IEEEtran}
\bibliography{refs}

\end{document}
